\scanpage{98}
Les fonctions $1,\unsure{f_\alpha|f}$ forment une base \inserted{sur $k$} de l'\unsure{ens|esp|ens.} des fonctions pol. de degré $\le2$, inv. par $G$.
% margin note written sideways
\margin{\uline{Conique circonscrite (canonique)} $\unsure{C(f)|C[f]|C(f_\alpha)}$ définie par $f=0$ (\unsure{indépndte|indépendte} de \deleted{\unsure{la}} normalisation). On vérifie que si $\alpha\ne0,1$ (i.e. cas du polyg. régulier à $n\ge5$ côtés) c'est la \uline{seule} conique circonscrite}
Parmi ces fonctions, $f$ est caractérisée par les \deleted{fonctions} conditions

a) $f_\alpha$ nulle sur $S$

b) La ``partie homogène'' \deleted{\unsure{d}}$f_0$ de $f$ satisfait
\[
f_0(s_i-s_{i-1})=1\qquad\forall i
\]
[NB
\[
\begin{aligned}
f_0(\deleted{\unsure{x,y}}X,Y)&=X^2+Y^2-\alpha XY\\
\text{donc}\qquad f_0(1,0)&=1\quad\big]
\end{aligned}
\]
\begin{center}\rule{1.4cm}{.35pt}\end{center}
7/ \deleted{Le discri} \uline{Discriminant de} $f_0$

Le discriminant de $f_0$ est
\[
\delta(f_0)=\deleted{\unsure{\alpha|2}}\;4-\alpha^2=(2-\alpha)(2+\alpha)
\]
qui est nul \unsure{sss|ssi} $\alpha=\pm2$ : la métrique définie par $f_0$ est non dég. sss on a $\alpha\ne\pm2$, i.e. sss on est dans le cas ``ordinaire''. On aura
\[
f_0(X,Y)=(Y-\unsure{\zeta|\rho}X)(Y\unsure{-|+}\zeta'X)
\]
où $\zeta,\zeta'$ sont les solutions de l'équation aux valeurs propres. Les ``\deleted{droites} \inserted{directions} isotropes''
% archive sheet number: 97
